% Part: normal-modal-logic
% Chapter: axioms-systems
% Section: normal-logics

\documentclass[../../../include/open-logic-section]{subfiles}

\begin{document}

\olfileid[fr]{nml}{prf}{nor}

\olsection{Logiques modales normales}

Tout ensemble de !!{formula}s modales ne se caractérise pas aisément
comme l'ensemble des formules !!{derivable}s à partir de certains
axiomes. Nous voulons que les logiques modales possèdent des
propriétés convenables. D'abord, tout ce que nous pouvons !!{derive}
en logique propositionnelle classique doit rester !!{derivable}, en
tenant compte du fait que les !!{formula}s peuvent maintenant
contenir aussi \iftag{prvBox}{$\Box$\iftag{prvDiamond}{
    et~}{}}{}\iftag{prvDiamond}{$\Diamond$}{}. Pour cela, nous
exigeons qu'une logique modale contienne toutes les instances
tautologiques et soit close par modus ponens.

\begin{defn}
  Une \emph{logique modale} est un ensemble~$\Sigma$ de
  !!{formula}s modales qui :
  \begin{enumerate}
  \item contient toutes les tautologies;
  \item est clos par substitution : si $!A \in \Sigma$ et si
    $!D_1$, \dots, $!D_n$ sont des !!{formula}s, alors
    \[
    \SSubst{!A}{\subst{!D_1}{p_1}, \dots, \subst{!D_n}{p_n}} \in \Sigma,
    \]
    \item est clos par \emph{modus ponens} : si $!A$ et $!A
      \lif !B \in \Sigma$, alors $!B \in \Sigma$.
  \end{enumerate}
\end{defn}

Pour utiliser la sémantique relationnelle des logiques modales, nous
devons en outre inclure toutes les !!{formula}s valides dans tous les
modèles modaux. Cette exigence est satisfaite dès que toutes les
instances de \Ax{K}\iftag{prvDiamond}{ et de~\Dual{}}{} sont
!!{derivable}s et que, si une !!{formula}~$!A$ est !!{derivable},
$\Box !A$ l'est également. Une logique modale qui remplit ces
conditions est dite \emph{normale}. (Il existe aussi des logiques
modales non normales, mais les modèles relationnels usuels ne leur
conviennent pas.)

\begin{defn}
  Une logique modale~$\Sigma$ est \emph{normale} si elle contient
  \begin{align*}
    \tag{\Ax{K}} & \Box(p \lif q) \lif (\Box p \lif \Box q),
    \iftag{prvDiamond}{\\
      \tag{\Dual} & \Diamond p \liff \lnot\Box\lnot p}{}
  \end{align*}
  et si elle est close par \emph{nécessitation} : si
  $!A \in \Sigma$, alors $\Box !A \in \Sigma$.
\end{defn}

Une implication tautologique transmet la \emph{vérité en un monde}
de son antécédent à son conséquent. La règle \Nec{} préserve la
\emph{vérité dans un modèle}, puisque si une formule est vraie en
tout monde, elle est vraie en tous les successeurs de chaque monde.
Elle préserve donc la validité dans un cadre et dans une classe de
cadres. Pour voir ce qui peut se passer en un monde donné, prenons
un monde~$w$ où $p$ est vraie et un successeur de~$w$ où $p$ est
fausse. En~$w$, $p$ est vraie et $\Box p$ est fausse.

\begin{prop}\ollabel{prop:rk}
  Toute logique modale normale est close par la règle \RK,
  \begin{prooftree}
    \AxiomC{$!A_1 \lif (!A_2 \lif \cdots (!A_{n-1} \lif !A_n)\cdots)$}
    \RightLabel{\RK}
    \UnaryInfC{$\Box!A_1 \lif (\Box!A_2 \lif \cdots (\Box!A_{n-1}
      \lif \Box!A_n)\cdots).$}
  \end{prooftree}
\end{prop}

\begin{proof}
  Raisonnons par récurrence sur~$n$. Si $n = 1$, la règle n'est autre
  que \Nec, et toute logique modale normale est close par \Nec.

  Supposons le résultat vrai pour $n-1$ et montrons-le pour~$n$.

  Supposons que
  \begin{align*}
  & !A_1 \lif (!A_2 \lif \cdots (!A_{n-1} \lif !A_n)\cdots) \in \Sigma
  \intertext{Par hypothèse de récurrence, nous avons}
  & \Box!A_1 \lif (\Box!A_2 \lif \cdots \Box(!A_{n-1} \lif !A_n)\cdots)
  \in \Sigma
  \intertext{Comme $\Sigma$ est une logique modale normale, elle
    contient toutes les instances de~$\Ax{K}$, en particulier}
  & \Box(!A_{n-1} \lif !A_n) \lif (\Box!A_{n-1} \lif \Box!A_n) \in \Sigma
  \intertext{Par modus ponens et à l'aide d'instances tautologiques
    appropriées, nous obtenons}
  & \Box!A_1 \lif (\Box!A_2 \lif \cdots (\Box!A_{n-1}
  \lif \Box!A_n)\cdots) \in \Sigma. 
  \end{align*}
\end{proof}

\begin{prop}\ollabel{prop:notDiamondBot}
  Toute logique modale normale~$\Sigma$ contient
  $\lnot\Diamond\lfalse$.
\end{prop}

\begin{prob}
  Démontrer la \olref[nml][prf][nor]{prop:notDiamondBot}.
\end{prob}

\begin{prop}
  Soient $!A_1$, \dots, $!A_n$ des !!{formula}s. Il existe une plus
  petite logique modale \emph{normale}~$\Sigma$ contenant toutes les
  instances de $!A_1$, \dots, $!A_n$.
\end{prop}

\begin{proof}
  Étant donnés $!A_1$, \dots, $!A_n$, définissons $\Sigma$ comme
  l'intersection de toutes les logiques modales normales contenant
  toutes les instances de $!A_1$, \dots, $!A_n$. Cette famille est
  non vide :
  $\Frm[L]$, l'ensemble de toutes les !!{formula}s, est une telle
  logique modale. Son intersection est encore une logique modale
  normale et contient les instances requises.
\end{proof}

\begin{defn}
La plus petite logique modale normale contenant $!A_1$, \dots,
$!A_n$ s'appelle un \emph{système modal} et se note
$\Log{K} !A_1 \dots !A_n$. La plus petite logique modale normale se
note~\Log{K}.
\end{defn}
\end{document}
